\documentclass[11pt]{article}

% Page layout and typography
\usepackage[margin=1in]{geometry}
\usepackage{setspace}
\usepackage{titlesec}
\usepackage{hyperref}
\usepackage{parskip}
\usepackage{booktabs}
\usepackage{array}
\usepackage[round]{natbib}

\setstretch{1.15}

% Heading style
\titleformat{\section}{\normalfont\bfseries}{\thesection}{1em}{}
\titleformat{\subsection}{\normalfont\bfseries}{\thesubsection}{1em}{}

% Link style
\hypersetup{
    colorlinks=true,
    urlcolor=blue,
    citecolor=black,
    linkcolor=black
}

\begin{document}

\begin{center}
{\LARGE \textbf{Paper Title Goes Here}}\\[0.5em]
{\large Optional Subtitle Goes Here}\\[0.8em]

Brandon Himpfen\\
Independent Researcher\\[0.5em]

Document type: Research report\\
Version: 1.0\\
Date: Month Day, Year\\
Source article: \url{https://www.example.com/source-article/}
\end{center}

\vspace{1em}

\section*{Abstract}

Write a concise abstract of 150--250 words. State the topic, purpose, method or approach, key findings, and intended use of the report. Keep the tone neutral, analytical, and research-oriented.

\textbf{Keywords:} keyword one; keyword two; keyword three; keyword four; keyword five

\section{Introduction}

Introduce the topic, problem, or research question. Explain why the issue matters and what this report contributes.

\section{Background}

Provide context, relevant prior work, market conditions, definitions, or conceptual framing.

\section{Methodology}

Describe how the report was prepared. For a synthesis paper, explain the sources reviewed, selection criteria, and analytical approach. For a framework paper, explain how the framework was developed.

\subsection{Data Sources}

Describe the sources used, such as public reports, datasets, industry publications, academic literature, or original analysis.

\subsection{Analytical Approach}

Describe how evidence was organized and interpreted.

\subsection{Limitations}

State limitations clearly. Examples include secondary-data constraints, inconsistent definitions, limited sample size, or lack of primary data.

\section{Findings or Framework}

Present the main analysis, framework, or findings.

\subsection{Dimension One}

Explain the first major component.

\subsection{Dimension Two}

Explain the second major component.

\subsection{Dimension Three}

Explain the third major component.

\begin{table}[h]
\centering
\caption{Example Evaluation Matrix}
\begin{tabular}{p{3.2cm} p{4.6cm} p{5cm}}
\toprule
\textbf{Dimension} & \textbf{Evaluation Question} & \textbf{Useful Evidence} \\
\midrule
Dimension One & What is being evaluated? & Evidence, indicators, metrics \\
Dimension Two & What is being evaluated? & Evidence, indicators, metrics \\
Dimension Three & What is being evaluated? & Evidence, indicators, metrics \\
\bottomrule
\end{tabular}
\end{table}

\section{Discussion}

Interpret the findings. Explain implications, tradeoffs, uncertainty, and areas where further research is needed.

\section{Conclusion}

Summarize the report's contribution and key takeaways. Keep the conclusion concise and avoid overclaiming.

\section*{Author Note}

This report is prepared as Version 1.0 for archival deposit and citation through Zenodo. Replace this note with any relevant context, source article link, repository link, or project note.

\section*{Recommended Citation}

Himpfen, B. (Year). \textit{Paper Title Goes Here} (Version 1.0) [Research report]. Zenodo. \url{https://doi.org/[DOI]}

\bibliographystyle{plainnat}
\bibliography{references}

\end{document}
